\documentclass[11pt]{article}

\usepackage[margin=1in]{geometry}
\usepackage{amsmath,amssymb,amsthm}
\usepackage{graphicx}
\usepackage[hidelinks]{hyperref}

\newtheorem{theorem}{Theorem}
\newtheorem{lemma}{Lemma}

\newcommand{\R}{\mathbb R}
\newcommand{\Jb}{\mathcal J_b^{1,\omega}(E)}
\newcommand{\Cb}{C_b^{1,\omega}(X)}
\newcommand{\Eop}{\mathcal E}

\title{A Literature-Implied Continuous Selection for Bounded \(C^{1,\omega}\) Whitney Jets}
\author{Agent lane 8 note for arXiv:1912.13317}
\date{2026-06-19}

\begin{document}
\maketitle

\section*{Status}

This is a literature-implied answer to the weaker continuity alternative in
Remark 4.3(2) of Azagra--Mudarra \cite{AzagraMudarra}.  The decisive external
input is the classical Bartle--Graves theorem: a continuous surjective linear
map between Banach spaces admits a continuous nonlinear right inverse
\cite{BartleGraves, BenyaminiLindenstrauss}.  The observation is that the
bounded \(C^{1,\omega}\) trace problem in \cite{AzagraMudarra} is exactly such
a quotient-map situation.

\section*{Source Question}

The source question appears on page 26 of the arXiv PDF, in Remark 4.3(2)
of \cite{AzagraMudarra}.

\begin{center}
\includegraphics[width=0.96\textwidth]{figures/open_problem_crop.png}
\end{center}

In words, the remark asks whether the bounded extension results can be
improved to obtain an extension operator with stronger \(A\)-seminorm
continuity, or at least an operator continuous from
\[
(\Jb,\rho)\quad\hbox{into}\quad(\Cb,\|\cdot\|_{1,\omega,b}).
\]

\section*{Identification}

The paper already proves, in Theorem 4.1, that every bounded Whitney jet
\((f,G)\in \Jb\) admits an extension \(F\in\Cb\), with
\[
\|F\|_{1,\omega,b}\le C\,\rho(f,G).
\]
The remaining point is not another extension formula.  It is a selection
problem for the linear trace operator.  Once the trace space is viewed with
the norm \(\rho\), Theorem 4.1 says precisely that this trace operator is
onto and has quotient norm equivalent to \(\rho\).  Bartle--Graves then gives
a continuous nonlinear choice of preimage.

\section*{Theorem}

\begin{theorem}
Let \(X\) be a Hilbert space, let \(E\subset X\), and let \(\omega\) be under
the hypotheses of Azagra--Mudarra's bounded extension theorem.  Then there is
a nonlinear extension operator
\[
S:(\Jb,\rho)\longrightarrow(\Cb,\|\cdot\|_{1,\omega,b})
\]
such that
\[
\bigl(S(f,G),\nabla S(f,G)\bigr)|_E=(f,G)
\]
for every \((f,G)\in\Jb\), and \(S\) is continuous.

Consequently, if \(\rho((f_n,G_n)-(f,G))\to0\), then
\[
\|S(f_n,G_n)-S(f,G)\|_{1,\omega,b}\to0.
\]
In particular,
\[
M_\omega\bigl(\nabla S(f_n,G_n)-\nabla S(f,G)\bigr)\to0
\]
and hence
\[
A\bigl(S(f_n,G_n)-S(f,G),
        \nabla S(f_n,G_n)-\nabla S(f,G)\bigr)\to0.
\]
\end{theorem}

\begin{proof}
Let
\[
B=(\Cb,\|\cdot\|_{1,\omega,b})
\]
and define the trace operator
\[
T:B\to \Jb,\qquad T(F)=(F|_E,\nabla F|_E).
\]
This map is linear.  It is bounded because
\[
\|F|_E\|_\infty\le \|F\|_\infty,\qquad
\|\nabla F|_E\|_\infty\le \|\nabla F\|_\infty,
\]
and, as recorded in the source paper before Corollary 1.6, every
\(C^{1,\omega}\) extension \(H\) of a jet satisfies
\[
A(f,G)\le M_\omega(\nabla H).
\]
Thus
\[
\rho(TF)\le \|F\|_{1,\omega,b}.
\]

Theorem 4.1 of \cite{AzagraMudarra} gives, for every \((f,G)\in\Jb\), an
extension \(F\in B\) with
\[
\|F\|_{1,\omega,b}\le C\rho(f,G).
\]
Therefore \(T\) is onto.  If \(q\) denotes the quotient norm induced by \(T\),
\[
q(f,G)=\inf\{\|F\|_{1,\omega,b}:T(F)=(f,G)\},
\]
then the two preceding estimates give
\[
\rho(f,G)\le q(f,G)\le C\rho(f,G).
\]
The quotient \(B/\ker T\) is a Banach space, and \(q\) is its norm under the
natural identification with the trace space.  Since \(q\) and \(\rho\) are
equivalent, \((\Jb,\rho)\) is Banach.

We may now apply the Bartle--Graves theorem to the bounded surjective linear
operator
\[
T:B\to(\Jb,\rho).
\]
It yields a continuous nonlinear map \(S:(\Jb,\rho)\to B\) such that
\(T\circ S=\mathrm{id}_{\Jb}\).  This identity is exactly the asserted
extension property, and continuity of \(S\) gives the norm convergence
claim.  The last displayed \(A\)-convergence follows from the source estimate
\(A(H,\nabla H)\le M_\omega(\nabla H)\), applied to
\[
H=S(f_n,G_n)-S(f,G).
\]
\end{proof}

\section*{Scope And Limitations}

The result answers the ``or at least'' part of Remark 4.3(2).  It also gives
the natural \(A\)-seminorm convergence for the difference of selected
extensions.  It does not show that Azagra--Mudarra's explicit
paraconvex-envelope operator has this continuity, nor does it give a new
closed formula for a continuous selector.  The proof is nonconstructive
through Bartle--Graves.

The displayed formula in the source remark appears to contain a typographical
mismatch in the gradient term.  The conclusion above is stated in the natural
form forced by convergence in \(\|\cdot\|_{1,\omega,b}\), namely for the
gradient of the difference of the selected extensions.

\section*{Search Notes}

Local run indexes were searched for `1912.13317', `C1omega', `1-jets',
`Whitney', `continuous extension operator', and `Bartle-Graves'.  Nearby
Whitney packets for arXiv:2305.19995 and arXiv:2403.14317 concern different
questions.  Web searches on 2026-06-19 found the source arXiv page and
general Bartle--Graves references, but no paper explicitly identifying the
above implication for Remark 4.3(2).  Accordingly this packet is recorded as
a literature-implied answer, not as an original construction.

\begin{thebibliography}{9}

\bibitem{AzagraMudarra}
D. Azagra and C. Mudarra,
\emph{\(C^{1,\omega}\) extension formulas for \(1\)-jets on Hilbert spaces},
arXiv:1912.13317.

\bibitem{BartleGraves}
R. G. Bartle and L. M. Graves,
\emph{Mappings between function spaces},
Transactions of the American Mathematical Society 72 (1952), 400--413.

\bibitem{BenyaminiLindenstrauss}
Y. Benyamini and J. Lindenstrauss,
\emph{Geometric Nonlinear Functional Analysis. Vol. 1},
American Mathematical Society Colloquium Publications 48, 2000.

\end{thebibliography}

\end{document}
